% Vietnamese translation for OI Public Library
% Source-File: IOI中国国家候选队论文/2005/杨俊/杨俊.doc
\documentclass[11pt]{article}
\usepackage{oipl-vn}

\newcommand{\Oh}{\mathcal{O}}
\newcommand{\Ans}{\mathrm{Ans}}
\newcommand{\true}{\mathrm{true}}
\newcommand{\false}{\mathrm{false}}

\title{Bàn sơ về ứng dụng của chiến lược nhị phân}
\author{Yang Jun}
\date{2005}

\begin{document}
\maketitle

\origtitle{Bàn sơ về ứng dụng của chiến lược nhị phân}
\sourcepath{IOI China candidate team papers / 2005 / Yang Jun.doc}

\renewcommand{\abstractname}{Tóm tắt}
\begin{abstract}
Bài viết tập trung thảo luận ba kiểu bài toán nhị phân khác nhau, với mục đích
giúp người đọc hiểu sâu hơn về phương pháp nhị phân. Đối tượng được xét đi từ
dãy có thứ tự thông thường, đến dãy có thứ tự đã suy biến, rồi cuối cùng là
dãy vô thứ tự. Trên thực tế, chúng cũng đại diện cho ba cách ứng dụng khác
nhau của chiến lược nhị phân.
\end{abstract}

\noindent\textbf{Từ khóa:} nhị phân; thứ tự; ứng dụng.

\section{Mở đầu}

``Nhị phân'' hẳn là một từ không còn xa lạ với mọi người. Nhị phân là một quy
tắc sàng lọc, bắt nguồn từ một ý tưởng rất đơn giản: trong trường hợp xấu
nhất, loại bỏ càng nhiều nhiễu càng tốt để tìm được mục tiêu nhanh nhất có
thể.

Hiệu quả của thuật toán nhị phân đến từ việc nó tận dụng đầy đủ thông tin,
cố gắng loại bỏ phần dư thừa và giảm những phép tính không cần thiết, từ đó
tối đa hóa hiệu suất của thuật toán. Thực ra, nhiều bài toán nhị phân có thể
được chứng minh bằng cây quyết định hoặc một số định lý khác rằng chúng đạt
tới cận dưới độ phức tạp của bài toán.

Dù bản thân tư tưởng nhị phân rất đơn giản, khả năng mở rộng và phạm vi ứng
dụng rộng của nó có lẽ vẫn vượt ngoài dự liệu ban đầu của chúng ta. Trong các
kỳ thi những năm gần đây, ta cũng thấy không ít ví dụ khiến người đọc bất ngờ
về cách dùng tư tưởng nhị phân. Dưới đây là ba kiểu ứng dụng khác nhau của tư
tưởng nhị phân mà tác giả tổng kết lại, hy vọng có ích cho người đọc.

\section{Dạng một: tìm kiếm nhị phân trên dãy có thứ tự thông thường}

Trước hết cần nói rõ: dãy có thứ tự ở đây không chỉ giới hạn ở dãy đã được sắp
từ nhỏ đến lớn hoặc từ lớn đến nhỏ như ta thường nói. Nó chỉ gồm hai ý: thứ
nhất, đó là một dãy một chiều; thứ hai, dãy ấy có thứ tự, tức là hai phần tử
bất kỳ trong dãy đều có thể so sánh được. Nói cách khác, nó có quan hệ thứ tự
toàn phần quen thuộc.

Tuy mọi người đều rất quen với tìm kiếm nhị phân, ta vẫn nhắc lại ngắn gọn quá
trình cài đặt tổng quát:
\begin{enumerate}
  \item Xác định phạm vi chứa phần tử cần tìm.
  \item Chọn một phần tử trong phạm vi đó làm mốc.
  \item So sánh khóa của phần tử cần tìm với khóa của phần tử mốc, rồi xác
  định phạm vi mới, chính xác hơn, của phần tử cần tìm.
  \item Nếu phạm vi mới đủ chính xác thì xuất kết quả; nếu không, quay lại
  bước 2.
\end{enumerate}

\subsection{Bài toán thống kê thứ tự}

\paragraph{Mô tả bài toán.}
Cho một tập $S$ gồm $n$ số đôi một khác nhau. Hãy tìm phần tử nhỏ thứ $i$ trong
tập đó.

\paragraph{Phân tích.}
Hẳn mọi người đều quen thuộc với bài toán này. Hãy nhắc lại cách áp dụng tìm
kiếm nhị phân vào bài toán:
\begin{enumerate}
  \item Xác định rằng phần tử cần tìm nằm trong $S$.
  \item Chọn ngẫu nhiên một trong $n$ phần tử, ký hiệu là $x$, và dùng $x$ làm
  mốc.
  \item Gọi $p$ là số phần tử trong $S$ nhỏ hơn $x$.
  \item Nếu $i<p$, phần tử cần tìm nhỏ hơn $x$, nên ta thu hẹp phạm vi về các
  phần tử của $S$ nhỏ hơn $x$ và tìm phần tử thứ $i$ trong phạm vi đó.
  \item Nếu $i>p$, phần tử cần tìm lớn hơn $x$, nên ta thu hẹp phạm vi về các
  phần tử của $S$ lớn hơn $x$ và tìm phần tử thứ $i-p-1$ trong phạm vi đó.
  \item Nếu $i=p$, phần tử nhỏ thứ $i$ chính là $x$.
  \item Nếu đã tìm được $x$ thì xuất kết quả; nếu không, quay lại bước 2.
\end{enumerate}
Vì $x$ được chọn ngẫu nhiên, một phân tích xác suất đơn giản cho thấy độ phức
tạp kỳ vọng của thuật toán là $\Oh(n)$.

\paragraph{Tiểu kết.}
Ví dụ này nhằm nhắc mọi người hai điểm.

Thứ nhất, đừng mặc nhiên cho rằng tìm kiếm nhị phân nhất định phải liên quan
đến $\log n$. Bước 3 trong thuật toán, tức bước ``chia'' theo cách nói thông
thường, không đòi hỏi lần nào cũng phải thực hiện trong thời gian $\Oh(1)$.
Việc chia có thể dựa trên xử lý hiệu quả đối với dãy, chẳng hạn như phép phân
hoạch tập hợp tương tự quicksort trong ví dụ trên.

Thứ hai, phép ``chia'' của thuật toán nhị phân không bắt buộc lần nào cũng
phải cân bằng. Đôi khi điều đó rất khó làm. Chỉ cần không lần nào cũng lệch
hẳn, ta đã có thể thu được một thuật toán hiệu quả; điều này cũng tạo cơ hội
cho việc dùng thuật toán ngẫu nhiên.

Trong những năm gần đây, sự xuất hiện của các bài tương tác cũng làm tìm kiếm
nhị phân thêm sinh động. Khá nhiều bài tìm kiếm nhị phân được đưa ra dưới dạng
bài tương tác. Chẳng hạn, ví dụ trên chỉ cần biến đổi một chút là thành một
bài tương tác độ khó trung bình ở IOI 2000. Hai bài gần như cùng khuôn: bạn
hỏi phần tử nhỏ thứ $i$, còn bài kia hỏi phần tử nhỏ thứ $(N+1)/2$; cách giải
cũng bắt chước tương tự. Bạn chọn ngẫu nhiên $x$ rồi chia thành hai đoạn theo
quan hệ lớn nhỏ với $x$; bài kia chọn ngẫu nhiên $x,y$ rồi chia thành ba đoạn
theo quan hệ lớn nhỏ với $x,y$. Độ phức tạp kỳ vọng của bạn là $\Oh(n)$, còn
số câu hỏi kỳ vọng của bài kia là $\Oh(N)$. Chi tiết cụ thể không được trình
bày ở đây; bạn đọc quan tâm có thể tham khảo báo cáo lời giải của các anh chị
đi trước.\footnote{Bài ``Báo cáo lời giải độ khó trung bình'' của Gao Yue,
tuyển tập luận văn đội tuyển huấn luyện 2000--2002.}

\section{Dạng hai: liệt kê nhị phân trên dãy có thứ tự suy biến}

Chiến lược nhị phân không phải lúc nào cũng được áp dụng trên một dãy có thứ
tự hiển nhiên như trên. Nó có thể dựa vào một mối liên hệ tiềm ẩn nào đó của
bài toán, rồi áp dụng lên một số dãy có thứ tự suy biến ngầm định. So với tìm
kiếm nhị phân đã giới thiệu ở trên, khác biệt lớn nhất là ở đây, khi quyết
định đệ quy sang phần nào, phép nhị phân không dùng phép so sánh trực tiếp.

\subsection{Giải quần vợt chuyên nghiệp BTP}

\paragraph{Mô tả bài toán.}
Bài này đến từ USACO contest tháng 12 năm 2002. Có $N$ con bò, trong đó
$N=2^K$, đều là cao thủ quần vợt. Mỗi con bò có một hạng BTP, đúng bằng một
số từ $1$ đến $N$. Tuần sau sẽ có một giải đấu loại trực tiếp: $N$ con bò được
chia thành $N/2$ cặp, mỗi cặp đấu để chọn ra $N/2$ người thắng; rồi $N/2$
người thắng lại được chia thành $N/4$ cặp, và cứ thế cho đến khi còn lại một
con bò vô địch.

Thi đấu vừa xét thực lực vừa xét vận may. Nếu hạng BTP của hai con bò chênh
lệch hơn số nguyên cho trước $K$, con bò có hạng tốt hơn luôn thắng con có
hạng kém hơn; ngược lại, hai bên đều có khả năng thắng. Khán giả muốn biết
con bò có hạng kém nhất trong số tất cả các con bò có thể trở thành vô địch
là con nào, đồng thời yêu cầu liệt kê một cách xếp lịch thi đấu khiến con bò
đó thắng. Giới hạn là $N\le 65536$.

\paragraph{Phân tích.}
Vì $N$ rất lớn, ta đoán rằng có thể giải bằng tham lam.

Ta muốn các tuyển thủ hạng tốt hơn luôn ``bất ngờ'' thua các tuyển thủ hạng
kém hơn. Vì vậy ta cho $1$ thua $K+1$, $2$ thua $K+2$, v.v. Trong mỗi ván của
mỗi vòng, luôn chọn tuyển thủ có hạng tốt nhất trong số chưa thi đấu, rồi cho
thua một tuyển thủ có hạng kém nhất có thể.

Nhưng rất dễ tìm phản ví dụ. Khi $N=8, K=2$, cách tham lam trên cho kết quả
là $4$ ở sơ đồ BTP-1, trong đó $X(Y$ nghĩa là $X$ thắng $Y$:
\[
\begin{array}{cccc}
3(1 & 4(2 & 7(5 & 8(6\\
4(3 &     & 8(7\\
4(8
\end{array}
\]
Trong khi đó, nghiệm tối ưu là $6$ ở sơ đồ BTP-2:
\[
\begin{array}{cccc}
6(7 & 5(3 & 4(8 & 2(1\\
6(5 &     & 4(2\\
6(4
\end{array}
\]
Nguyên nhân là ta không biết nghiệm tối ưu bằng bao nhiêu, nên đang tham lam
một cách mù quáng. Trên thực tế, nghiệm tối ưu $6$ đã bị ta loại ngay ở vòng
đầu, vì thế tất nhiên không thể thu được nghiệm tối ưu.

Muốn biết nghiệm tối ưu? Hãy liệt kê. Đồng thời, xét một kết luận rất hiển
nhiên: nếu tuyển thủ hạng $X$ có thể thắng chung cuộc, thì mọi tuyển thủ có
hạng tốt hơn $X$ cũng có thể thắng chung cuộc.

Dù kết luận có vẻ hiển nhiên, để chứng minh nó vẫn cần suy nghĩ một chút.
Giả sử tuyển thủ hạng $X$ thắng chung cuộc. Ta sửa cục bộ cách xếp lịch ấy để
xây dựng một cách xếp lịch mới khiến một tuyển thủ bất kỳ hạng $Y<X$ thắng.
Trong lịch mà $X$ thắng chung cuộc, giả sử $Y$ bị $Z$ đánh bại. Nếu $Z\ne X$,
do $X>Y$, ta biết $Z$ cũng có thể đánh bại $X$, và trong cùng vòng ấy cũng như
các vòng sau, mọi đối thủ mà $X$ đánh bại đều có thể bị $Y$ đánh bại. Vì vậy
ta chỉ cần cho $Z$ đánh bại $X$ ở vòng này, rồi đổi mọi lần xuất hiện của $X$
trong các trận sau thành $Y$. Nếu $Z=X$, do $X>Y$, ta cho $Y$ đánh bại $X$,
rồi cũng đổi mọi trận sau của $X$ thành $Y$; cuối cùng người thắng vẫn là
$Y$.

Nhờ đó, ta có thể liệt kê nhị phân giá trị $X$ của người thắng chung cuộc,
nâng cao hiệu quả thuật toán đáng kể.

Bây giờ vấn đề là: nếu đã biết người thắng chung cuộc là $X$, ta có thể nhanh
chóng xây dựng một cách xếp lịch, hoặc chứng minh rằng không tồn tại, hay
không? Có. Ta vẫn dùng tham lam, nhưng vì đã biết ai thắng cuối cùng nên ta
làm ngược. Ở mỗi vòng, ta để các tuyển thủ đã có mặt đánh bại tuyển thủ có
hạng tốt nhất trong số chưa xuất hiện. Lịch đấu sinh ra bởi cách này chính là
kiểu được mô tả trong BTP-2. Cách tìm tuyển thủ chưa xuất hiện có hạng tốt
nhất có thể dùng cây nhị phân tìm kiếm tĩnh; phần này không trình bày thêm,
bạn đọc có thể tự suy nghĩ.

Có thể chứng minh cách tham lam này đúng; phương pháp chứng minh tương tự
chứng minh phía trên nên không lặp lại. Toàn bộ bài toán có thể hoàn thành
trong thời gian $\Oh(N\log_2 N)$.

\paragraph{Tiểu kết.}
Đối với dạng bài cần liệt kê nhị phân này, bản chất của thuật toán là tìm kiếm
nhị phân trên một dãy có thứ tự suy biến ngầm định, chỉ có hai giá trị $0$ và
$1$. Trong ví dụ trên, khi $X<\Ans$ thì $A[X]=0$, còn khi $X\ge\Ans$ thì
$A[X]=1$. Điều ta cần tìm chính là điểm phân giới giữa hai loại giá trị ấy. Có
điểm phân giới thì có giá trị tối ưu, và nhờ đó có lời giải cho bài toán ban
đầu.

\subsection{Phân phát huy chương}

\paragraph{Mô tả bài toán.}
Bài này đến từ ACM/ICPC CERC 2004. Có một bức tường hình vòng tròn; trên
tường có một số tháp, mỗi tháp có một lính gác. Bây giờ cần phát cho mỗi lính
gác một số huy chương. Lính gác thứ $i$ cần $P[i]$ huy chương. Các huy chương
của cùng một lính gác phải thuộc các loại khác nhau, và hai lính gác kề nhau
không được có bất kỳ huy chương nào cùng loại. Hỏi ít nhất cần chuẩn bị bao
nhiêu loại huy chương khác nhau. Giới hạn là $2\le n\le 10000$ và
$1\le P[i]\le 100000$.

\paragraph{Phân tích.}
Nếu bức tường không phải hình vòng tròn, ta rất dễ giải bằng tham lam: mỗi lần
phát cho lính gác nhiều nhất có thể từ các huy chương đã chuẩn bị; nếu không
đủ thì chuẩn bị thêm một số loại để thỏa mãn nhu cầu. Nhưng bây giờ bức tường
là vòng tròn, lính gác cuối cùng và lính gác đầu tiên cũng không được trùng
loại huy chương, và đây là điểm khó của bài toán.

Vì vậy ta thử đưa điểm khó này vào trạng thái và dùng quy hoạch động
\footnote{Nói chính xác hơn ở đây là truy hồi.} để giải, vì quy hoạch động là
một phương pháp thường dùng cho các bài toán tối ưu.

Gọi $a[i,j]$ là số loại huy chương tối thiểu cần thêm, ngoài $P[1]$ loại huy
chương đã cấp cho lính gác đầu tiên, sau khi đã sắp xếp xong $i$ lính gác đầu
và lính gác thứ $i$ có $j$ huy chương cùng loại với lính gác đầu tiên.

Giả sử lính gác thứ $i-1$ có $k$ huy chương cùng loại với lính gác đầu tiên.
Vì $k$ huy chương ấy và $j$ huy chương của lính gác thứ $i$ phải đôi một khác
loại, ta có ngay $j+k\le P[1]$. Mặt khác, $P[i-1]-k$ huy chương còn lại của
lính gác thứ $i-1$ cũng phải khác với $P[i]-j$ huy chương còn lại của lính
gác thứ $i$. Gọi số loại huy chương cần tăng thêm là $X$, ta có
\[
  a[i-1,k]+X \ge (P[i]-j)+(P[i-1]-k),
\]
suy ra
\[
  X \ge (P[i]-j)+(P[i-1]-k)-a[i-1,k].
\]
Do đó có thể viết chuyển trạng thái dưới dạng
\[
  a[i,j]=
  \min_{\substack{k\\ j+k\le P[1]}}
  \max\{a[i-1,k],\, (P[i]-j)+(P[i-1]-k)\}.
\]
Làm trực tiếp như vậy có độ phức tạp tới $\Oh(nP_{\max}^2)$. Dù có tối ưu,
ta cũng khó thu được kết quả thỏa đáng, nên cần tìm lối khác.

Xét tính chất sau: nếu với $P[1]+X$ loại huy chương tồn tại một phương án
phân phát thỏa yêu cầu, thì với $P[1]+X+1$ loại cũng chắc chắn có phương án
khả thi; cùng lắm ta không dùng một loại nào đó. Tính chất này rất quan trọng,
vì nhờ nó, ta có thể liệt kê nhị phân $X$ rồi lần lượt kiểm tra tính khả thi
để thu được kết quả.

Vậy bài toán ban đầu được chuyển thành: nếu đã biết tổng cộng có $P[1]+X$ loại
huy chương, liệu ta có thể nhanh chóng kiểm tra có tồn tại một cách phân phát
thỏa yêu cầu hay không? Câu trả lời là có.

Phương pháp vẫn là quy hoạch động, và cách biểu diễn trạng thái cũng tương tự.
Gọi $a[i,j]$ cho biết, khi tổng số loại huy chương là $P[1]+X$, có thể sắp xếp
xong $i$ lính gác đầu và làm cho lính gác thứ $i$ có $j$ huy chương cùng loại
với lính gác đầu tiên hay không. Khi đó chuyển trạng thái trở thành:
\[
  a[i,j]=\true
  \Longleftrightarrow
  \exists k:\,
  \begin{cases}
    a[i-1,k]=\true,\\
    j+k\le P[1],\\
    (P[i]-j)+(P[i-1]-k)\le X.
  \end{cases}
\]
Nhìn bề ngoài, chuyển trạng thái vẫn khá rườm rà, vì $k$ có rất nhiều ràng
buộc. Nhưng quan sát kỹ, hai điều kiện thứ hai và thứ ba hợp lại thành
\[
  P[i-1]+P[i]-X-j \le k \le P[1]-j.
\]
Với $i$ cố định, miền giá trị của $k$ trên trục số là một đoạn có độ dài xác
định, và vị trí của đoạn do $j$ quyết định. Ràng buộc thứ nhất ban đầu chỉ là
giới hạn thêm phạm vi cho đoạn này, loại bỏ phần dư thừa vô nghĩa.

Trạng thái ban đầu chỉ có $a[1,0]=\true$, các trạng thái còn lại đều là
$\false$. Bằng quy nạp, rất dễ chứng minh rằng với mọi $i$, tập các $j$ sao
cho $a[i,j]$ là $\true$ nhất định là một đoạn liên tiếp. Vì vậy, với mỗi $i$,
tất cả các trạng thái $a[i,j]$ có thể được biểu diễn chỉ bằng hai đầu mút
$a[i].x1$ và $a[i].x2$.

Để $a[i,j]=\true$, khi và chỉ khi miền giá trị của $k$ giao với đoạn
$[a[i-1].x1,a[i-1].x2]$, tức là thỏa
\[
  P[1]-j \ge a[i-1].x1
  \quad\text{và}\quad
  P[i-1]+P[i]-X-j \le a[i-1].x2.
\]
Nói cách khác,
\[
  P[i-1]+P[i]-X-a[i-1].x2 \le j \le P[1]-a[i-1].x1.
\]
Từ đó, phương trình chuyển trạng thái trở thành
\[
  a[i].x1=\max(0,\ P[i-1]+P[i]-X-a[i-1].x2),
\]
\[
  a[i].x2=\min(P[i],\ P[1]-a[i-1].x1).
\]
Trạng thái ban đầu là $a[1].x1=a[1].x2=P[1]$. Số trạng thái là $\Oh(n)$, mỗi
chuyển trạng thái tốn $\Oh(1)$, có thể nói là rất hiệu quả.

Với phương pháp này, dựa trên phân tích trước đó, ta liệt kê nhị phân $X$ và
dùng cách kiểm tra hiệu quả nói trên, từ đó giải toàn bộ bài toán trong thời
gian $\Oh(n\log P_{\max})$.

\paragraph{Tiểu kết.}
Bài toán trên nhìn qua tưởng như điểm khó vẫn là quy hoạch động, còn liệt kê
nhị phân chỉ là một công cụ phụ. Nhưng trên thực tế, việc liệt kê trước $X$ đã
làm phương trình quy hoạch động đơn giản đi rất nhiều, nhờ đó bài toán mới
được giải quyết. Những bài tối ưu dùng tư tưởng liệt kê nhị phân như thế gần
đây khá thịnh hành, và rất dễ dẫn dụ thí sinh trực tiếp dùng quy hoạch động
hoặc tham lam rồi đi vào ngõ cụt.

Vì vậy, về sau khi xét các bài toán tối ưu, ta nên chú ý xem bài toán có ẩn
chứa một dãy $01$ có thứ tự hay không, và liệu có thể dùng liệt kê nhị phân để
chuyển bài toán tối ưu thành bài toán khả thi hay không. Hãy ghi nhớ: lùi một
bước, trời biển mở rộng.

\section{Dạng ba: tìm kiếm nhị phân trên dãy vô thứ tự}

Nếu bạn cho rằng chỉ trong dãy có thứ tự mới có thể dùng nhị phân, thì bạn đã
sai lớn rồi; dãy vô thứ tự vẫn có thể được tìm kiếm bằng nhị phân. Hãy xem bài
tương tác sau.

\subsection{Chuyến đi của người bán hàng}

\paragraph{Mô tả bài toán.}
Bài này đến từ JSOI. Là một người bán hàng, Mirko phải đi máy bay thăm $N$
thành phố, mỗi thành phố đúng một lần. Giữa mỗi cặp thành phố có đúng một
đường bay; máy bay luôn cất cánh vào đúng giờ và bay mất một giờ. Mỗi tuyến
bay luôn có máy bay bay qua lại liên tục giữa hai thành phố: nếu có một máy
bay lúc 5 giờ cất cánh từ thành phố $A$ đến thành phố $B$, thì 6 giờ nó đến
nơi, lập tức lại cất cánh, và 7 giờ quay về $A$, v.v.

Để bảo đảm hiệu quả, Mirko muốn sắp $N$ thành phố thành một dãy
$A_1,A_2,\ldots,A_N$ sao cho với mỗi $i$ ($1<i<n$), sau khi tới thành phố
$A_i$, Mirko có thể quảng cáo một giờ rồi lập tức rời $A_i$ để đi tới
$A_{i+1}$.

Đáng tiếc là Mirko không có lịch bay, nên anh phải gọi điện hỏi hãng hàng
không. Mỗi cuộc gọi, anh có thể hỏi tuyến giữa hai thành phố $A,B$ vào đúng
12 giờ trưa bay từ $A$ sang $B$ hay từ $B$ sang $A$. Vì Mirko không muốn tốn
quá nhiều tiền điện thoại, hãy giúp anh xác định một lộ trình bằng ít cuộc gọi
nhất có thể. Chú ý rằng thư viện tương tác có thể điều chỉnh tuyến theo các
câu hỏi của bạn để làm số cuộc gọi của bạn lớn nhất.

\paragraph{Phân tích.}
Bài toán khá dài, trước hết hãy phân tích xem rốt cuộc nó yêu cầu gì.

Giả sử 12 giờ trưa ta bay từ $A_1$ tới $A_2$, 1 giờ đến nơi, rồi quảng cáo một
giờ, 2 giờ rời $A_2$ để tới $A_3$, v.v. Rõ ràng, ta luôn rời một thành phố để
tới thành phố khác vào các giờ chẵn.

Vì thời gian bay khứ hồi đúng bằng 2 giờ và ở giữa không có thời gian nghỉ,
nên nếu chuyến 12 giờ bay từ $A$ tới $B$, thì vào mọi giờ chẵn, chuyến bay
cũng nhất định bay từ $A$ tới $B$.

Nhờ đó, ta có thể chuyển bài toán ban đầu thành bài toán tìm một đường
Hamilton trong một tournament,\footnote{Tức là một đồ thị có hướng gồm $N$
đỉnh, trong đó giữa mỗi cặp đỉnh có đúng một cung theo một trong hai chiều.}
và mục tiêu là dùng ít câu hỏi nhất có thể. Lưu ý đây là một phép chuyển hóa,
không phải tương đương hoàn toàn.

Để tránh hỏi một cách mù quáng, ta thử dùng phương pháp tăng dần để mở rộng
dãy. Trước hết hỏi tùy ý hai thành phố; cho tiện, chọn thành phố 1 và thành
phố 2. Không mất tính tổng quát, giả sử có cung từ 1 đến 2, ta thu được một
đường độ dài 1 là $1(2$.

Giả sử ta đã thiết kế được một đường gồm $i$ thành phố đầu, độ dài $i-1$:
$A_1(A_2(\cdots(A_i$. Ta muốn thêm thành phố $i+1$ vào để thành đường độ dài
$i$.

Trước hết hỏi có cung từ $i+1$ tới $A_1$ hay không. Nếu có, ta đổi đường thành
$i+1(A_1(A_2(\cdots(A_i$. Nếu không, ta hỏi tiếp có cung từ $A_i$ tới $i+1$
hay không. Nếu có, ta đổi đường thành $A_1(A_2(\cdots(A_i(i+1$. Nếu vẫn
không, điều đó nghĩa là không thể chèn ở hai đầu, nên $i+1$ chỉ có thể chèn
vào giữa.

Chú ý rằng cạnh giữa $i+1$ và $A_1$ khi ấy có chiều từ $A_1$ tới $i+1$. Nếu có
cung từ $i+1$ tới $A_2$, ta có thể chèn $i+1$ giữa $A_1$ và $A_2$; nếu $i+1$
vẫn không có cung tới $A_2$, tức là có cung từ $A_2$ tới $i+1$, ta hỏi tiếp
$A_3$, v.v. Liệu cách hỏi này có thể đi tới cuối cùng mà vẫn không chèn được
$i+1$ hay không? Không, vì $i+1$ có cung tới $A_i$.

Như vậy ta có một phương pháp hỏi, nhưng trong trường hợp xấu nhất, tổng số
câu hỏi có thể là $\Oh(N^2)$.

Với mỗi điểm $A_k$ ($1\le k\le i$), hoặc $A_k$ có cung tới $i+1$, hoặc $i+1$
có cung tới $A_k$. Hơn nữa $A_1$ có cung tới $i+1$, còn $i+1$ có cung tới
$A_i$. Do đó, ta có thể dùng tìm kiếm nhị phân để tìm $A_k$ sao cho $i+1$ có
thể được chèn giữa $A_k$ và $A_{k+1}$.

Giả sử ta đã biết $i+1$ có thể được chèn giữa $A_p$ và $A_r$. Ta hỏi chiều
cung giữa $i+1$ và $A_{(p+r)/2}$. Nếu $i+1$ có cung tới $A_{(p+r)/2}$, điều
đó cho biết $i+1$ có thể được chèn trong đoạn từ $A_p$ đến $A_{(p+r)/2}$.
Ngược lại, nó có thể được chèn trong đoạn từ $A_{(p+r)/2}$ đến $A_r$. Cứ nhị
phân như vậy cho đến khi $p+1=r$.

Nhờ đó, số câu hỏi cần dùng chỉ là $\Oh(n\log n)$.

\paragraph{Tiểu kết.}
Đối với dạng tìm kiếm nhị phân này, về bản chất ta đang tìm kiếm trên một dãy
$01$ vô thứ tự; đối tượng cần tìm là một xâu con đặc biệt ``01''. Có xâu con
này rồi, dùng phương pháp xây dựng, ta có thể chuyển kết quả thành một nghiệm
khả thi của bài toán ban đầu.

Khi hỏi cung giữa $i+1$ và $A_{(p+r)/2}$, nếu $i+1$ có cung tới
$A_{(p+r)/2}$, ta tìm cách chèn $i+1$ vào giữa $A_p$ và $A_{(p+r)/2}$. Thực ra
điều đó không có nghĩa là $i+1$ chắc chắn không thể chèn vào giữa
$A_{(p+r)/2}$ và $A_r$; chỉ là chưa chắc có thể. Còn $i+1$ nhất định có thể
chèn vào giữa $A_p$ và $A_{(p+r)/2}$. Từ đó có thể thấy, phương pháp tìm kiếm
nhị phân kiểu này thường áp dụng cho những bài toán có nhiều nghiệm khả thi,
nhưng cần xây dựng một nghiệm khả thi một cách hiệu quả.

\subsection{Mạch cổng NAND}

\paragraph{Mô tả bài toán.}
Bài này đến từ ACM/ICPC CERC 2001. Có một loại cổng mạch gọi là cổng NAND.
Mỗi cổng NAND có hai đầu vào, đầu ra là kết quả của phép NAND trên hai đầu
vào, tức $\operatorname{not}(A\operatorname{and}B)$. Cho một mạch không chu
trình gồm $N$ cổng NAND; tính không chu trình rất quan trọng đối với một mạch
logic. Tất cả $M$ đầu vào của mạch đều nối với cùng một đầu vào $x$, như sơ đồ
NAND-1. Hãy đặt một số đầu vào thành hằng số, dùng số lượng $x$ ít nhất để
hoàn thành cùng chức năng. NAND-2 là một mạch chỉ dùng một đầu vào $x$ nhưng
có thể thu được cùng kết quả.

\paragraph{Phân tích.}
Ý tưởng đầu tiên là dùng biểu thức: biểu diễn đầu ra của mỗi cổng mạch thành
kết quả của một phép tính hàm nào đó theo đầu vào, từ đó nối đầu vào và đầu ra
của toàn mạch. Nhưng phép NAND không thỏa tính kết hợp, nên không thể tính
toán biểu thức như vậy. Ta đành tìm một con đường khác.

Thực ra, ta rất dễ tính được đầu ra của mạch khi $x$ lần lượt bằng $0$ và
$1$, ký hiệu là $\mathit{ans0}$ và $\mathit{ans1}$. Nếu
$\mathit{ans0}=\mathit{ans1}$, rõ ràng ta không cần dùng đầu vào $x$ nào; chỉ
cần đặt toàn bộ đầu vào thành $0$ hoặc $1$ là thỏa yêu cầu. Nếu
$\mathit{ans0}\ne\mathit{ans1}$ thì sao? Không dùng đầu vào nào là điều không
thể. Vậy chỉ dùng một đầu vào thì sao?

Xét rằng, ngay cả khi chỉ có một $x$, số dãy đầu vào ta có thể chọn vẫn rất
lớn; chỉ cần có một dãy thỏa yêu cầu là đủ. Vì vậy ta có thể mạnh dạn phỏng
đoán rằng nhiều nhất chỉ cần một đầu vào chưa biết $x$ là có thể thỏa yêu
cầu.

Chỉ dùng một $x$ nghĩa là khi $x=0$, đầu ra của mạch đúng bằng $\mathit{ans0}$;
khi $x=1$, đầu ra của mạch đúng bằng $\mathit{ans1}$. Ta chỉ thay đổi giá trị
của một cổng đầu vào, và thay từ $0$ thành $1$.

Từ đó, ta tưởng tượng một dãy đầu vào toàn $0$. Sau một số lần sửa đổi như
trên, cuối cùng nó sẽ trở thành một dãy đầu vào toàn $1$. Mặt khác, lúc bắt
đầu đầu ra của mạch là $\mathit{ans0}$, còn cuối cùng là $\mathit{ans1}$. Vì
vậy, chắc chắn tồn tại một thời điểm mà khi một đầu vào đổi từ $0$ thành $1$,
đầu ra vừa đúng đổi từ $\mathit{ans0}$ sang $\mathit{ans1}$.

Điều này chứng minh phỏng đoán vừa nêu: nhiều nhất chỉ cần một $x$. Đồng thời,
thuật toán cũng hiện ra. Giả tưởng rằng ta xây dựng $M+1$ bộ đầu vào khác
nhau như sau:
\[
\begin{array}{c@{\quad}l}
0: & 0\ 0\ 0\ 0\ 0\ 0\ \cdots\ 0,\\
1: & 1\ 0\ 0\ 0\ 0\ 0\ \cdots\ 0,\\
2: & 1\ 1\ 0\ 0\ 0\ 0\ \cdots\ 0,\\
\vdots & \vdots\\
M: & 1\ 1\ 1\ 1\ 1\ 1\ \cdots\ 1.
\end{array}
\]
Ta cần tìm hai bộ đầu vào kề nhau sao cho bộ trước cho đầu ra
$\mathit{ans0}$, còn bộ sau cho đầu ra $\mathit{ans1}$.

Việc này vẫn có thể hoàn thành hiệu quả bằng tìm kiếm nhị phân. Biết rằng bộ
thứ $p$ cho đầu ra $\mathit{ans0}$ và bộ thứ $r$ cho đầu ra $\mathit{ans1}$,
ta tính đầu ra $\mathit{ans'}$ của bộ thứ $(p+r)/2$. Nếu
$\mathit{ans'}=\mathit{ans0}$ thì lấy $r=(p+r)/2$; ngược lại lấy
$p=(p+r)/2$, cho đến khi $p+1=r$.

Như vậy bài toán được giải quyết gọn trong thời gian $\Oh(N\log M)$.

\paragraph{Tiểu kết.}
Có thể thấy, điểm khó của dạng tìm kiếm nhị phân này nằm ở xây dựng. Làm thế
nào đưa một bài toán tưởng như hoàn toàn không liên quan vào phạm vi ứng dụng
của tìm kiếm nhị phân, rồi dùng một thay đổi nhỏ từ $0$ sang $1$ để xây dựng
một lời giải cho bài toán ban đầu? Phương pháp xây dựng không có khuôn mẫu hay
quy luật thống nhất; chỉ có thể phân tích cụ thể từng bài. Điều ta có thể làm
là tiếp xúc nhiều hơn, suy nghĩ nhiều hơn và tích lũy nhiều hơn trong quá
trình luyện tập.

\section{Tổng kết}

Bài viết này chỉ giới thiệu ngắn gọn vài ví dụ; việc bao quát toàn bộ, thậm
chí phần lớn, các ứng dụng của chiến lược nhị phân là rất khó. Điều tôi muốn
chỉ ra là: dù tư tưởng nhị phân đơn giản, nội dung của nó vẫn rất phong phú.
Ta không nên để tư tưởng nhị phân chỉ dừng lại ở phép tìm kiếm nhị phân cơ bản
nhất trên mảng có thứ tự thông thường, mà nên mở rộng nó sang các ứng dụng
rộng hơn.

Chiến lược nhị phân thường kết hợp với một số thuật toán khác, nhờ đó che giấu
bản thân và thoát khỏi tầm nhìn của chúng ta. Điều này đòi hỏi ta phải có nền
tảng cơ bản vững chắc, kinh nghiệm giải bài phong phú, những phỏng đoán táo
bạo nhưng hợp lý, cùng tư duy sáng tạo linh hoạt, thì mới có thể ``lấy bất
biến ứng vạn biến''.

\section*{Tài liệu tham khảo}

\begin{enumerate}
  \item Wu Wenhu và Wang Jiande, \textit{Phân tích và thiết kế chương trình
  cho các thuật toán thực dụng}.
  \item Liu Rujia và Huang Liang, \textit{Nghệ thuật thuật toán và thi đấu tin
  học}.
\end{enumerate}

\end{document}
